\chapter{Tachycardia}

\section*{Definition}
Tachycardia is when your heart beats faster than 100 times per minute while at rest. It can start in different parts of your heart and may come and go or be constant.

\section*{What Happens in Your Body}
A rapid heart rate can happen when your heart's electrical system fires too quickly or creates short-circuit loops. Sinus tachycardia is a normal response to exercise or fever, while other types like atrial fibrillation need medical treatment.

\section*{Causes}
\begin{itemize}
  \item Sinus rapid heart rate (fever, dehydration, anxiety, exercise)
  \item Atrial fibrillation or atrial flutter
  \item Supraventricular rapid heart rate (SVT, AVNRT, WPW)
  \item Ventricular rapid heart rate (VT)
  \item Hyperthyroidism
  \item Anemia
  \item Pulmonary embolism
  \item Heart failure or cardiomyopathy
  \item Electrolyte imbalances
  \item Medication or substance effects (caffeine, alcohol, stimulants, thyroid medication)
\end{itemize}

\section*{When to See a Doctor}
See your doctor if:
\begin{itemize}
  \item Symptoms are severe or getting worse over time
  \item Resting rapid heart rate with chest pain, shortness of breath, fainting, or dizziness (possible cardiac arrhythmia)
  \item You have other concerning symptoms such as fever, weight loss, or bleeding
\end{itemize}

\section*{Related Symptoms}
\begin{itemize}
  \item Palpitations
  \item Shortness of breath
  \item Dizziness
  \item Chest pain
  \item Fatigue
\end{itemize}
\textbf{Important:} If this symptom persists, your doctor will check for serious underlying causes.

\section*{Self-Care / Home Management}
\begin{itemize}
  \item For acute palpitations without chest pain or syncope, sit down, breathe slowly and deeply, and avoid caffeine, nicotine, and energy drinks.
  \item Perform vagal manoeuvres (bearing down as if having a bowel movement or coughing forcefully) for potentially benign supraventricular tachycardia to abort the episode.
  \item Maintain adequate hydration (1.5--2 L water daily) and electrolyte balance; avoid excessive alcohol, stimulants, and non-prescribed thyroid hormones.
  \item Keep a symptom diary noting heart rate, triggers, duration, and associated symptoms to share with the evaluating physician.
\end{itemize}

\section*{How Your Doctor Will Evaluate You}
\begin{itemize}
  \item Twelve-lead ECG to document rhythm, rate, QRS width, and signs of ischaemia or pre-excitation (delta wave in WPW); an ECG during symptoms is the most valuable diagnostic tool.
  \item Laboratory workup: TSH, CBC, electrolytes (potassium, magnesium, calcium), cardiac troponin if chest pain or cardiac history, and drug screen if stimulant use is suspected.
  \item Ambulatory cardiac monitoring: 24-hour Holter monitor for daily symptoms, event recorder for intermittent symptoms (seven to 30 days), or mobile cardiac telemetry for prolonged monitoring.
  \item Echocardiogram to evaluate structural heart disease: valvular abnormalities, left ventricular function, wall motion abnormalities, and pericardial effusion.
\end{itemize}

\section*{Treatment Options}
\begin{itemize}
  \item Sinus tachycardia: treat the underlying cause (fever, dehydration, anaemia, hyperthyroidism) rather than the heart rate itself; beta-blockers for symptomatic control.
  \item Atrial fibrillation or flutter: rate control (beta-blocker or calcium channel blocker) or rhythm control with antiarrhythmics; anticoagulate per CHA2DS2-VASc score to prevent stroke.
  \item Supraventricular tachycardia: acute management with vagal manoeuvres or IV adenosine 6--12 mg; catheter ablation is curative for most SVT (AVNRT, AVRT/WPW).
  \item Ventricular tachycardia: treat as a medical emergency with IV amiodarone 150 mg over ten minutes or synchronised cardioversion; implantable cardioverter-defibrillator placement for secondary prevention.
\end{itemize}

\section*{References}
\begin{thebibliography}{99}
\bibitem{tachycardia:esc-tachy} Zeppenfeld K, et al. ``2023 ESC Guidelines for the management of ventricular arrhythmias.'' \emph{Eur Heart J}. 2023;44(40):3965-4048.
\bibitem{tachycardia:uptodate-tachy} Page RL, et al. Approach to the patient with tachycardia. In: UpToDate, Post TW (Ed), UpToDate, Waltham, MA; 2025.
\bibitem{tachycardia:nejm-svt} Brugada J, et al. ``Supraventricular Tachycardia.'' \emph{N Engl J Med}. 2023;389(7):629-640.
\bibitem{tachycardia:jam-tachy} Al-Khatib SM, et al. ``Management of atrial fibrillation.'' \emph{JAMA}. 2024;331(3):236-247.
\bibitem{tachycardia:harrison-tachy} Jameson JL, et al. \emph{Harrison\'s Principles of Internal Medicine}. 21st ed. McGraw-Hill; 2022. Chapter 244: Approach to Tachyarrhythmias.
\bibitem{tachycardia:cochrane-tachy} Chen LY, et al. ``Catheter ablation versus medical therapy for atrial fibrillation.'' \emph{Cochrane Database Syst Rev}. 2023;(10):CD013578.
\end{thebibliography}
